% SEGMENT OLP-0702-S01
% Part: proof-theory
% Chapter: sequent-calculus
% Section: proof-examples

\documentclass[../../../include/open-logic-section]{subfiles}

\begin{document}

\olfileid{pt}{seq}{ex}

\olsection{நிறுவல்களின் எடுத்துக்காட்டுகள்}

தொடரணிகளுக்கான !!{proof}s அமைக்கும்போது,
முடிவுத் தொடரணியிலிருந்து மேல்நோக்கிச் செல்வது
பொதுவாக வசதியானது. அதிலுள்ள ஓர்
!!a{formula} வைச் செயலுக்குரிய !!{formula} வாகத் தேர்ந்து,
அதற்குரிய விதியின் முன்கூற்றுகளைத் தொடரணியின்
மேலே எழுதி, அடிக்கோள்களை அடையும்வரை
இம்முறையைத் தொடரலாம். எடுத்துக்காட்டாக,
பின்வரும் தொடரணியை நிறுவ வேண்டும் என்க:
\[(!C \land !D) \lif !E \Sequent !C \lif (!D \lif !E)\]
$!C \lif (!D \lif !E)$ ஐப் பார்ப்போம்: இது
$!A \lif !B$ என்ற வடிவில் தொடரணியின்
வலப்புறத்தில் உள்ளது. முழுத் தொடரணியையும்
\Log{G1c} இன் $\RightR{\lif}$ விதியுடன்
ஒப்பிட்டால்,
\[\Axiom$ !A, \Gamma \fCenter \Delta, !B$
\RightLabel{\RightR{\lif}}
\UnaryInf$ \Gamma \fCenter \Delta, !A \lif !B $
\DisplayProof\]
இங்கு $\Gamma = \{(!C \land !D) \lif !E\}$,
$\Delta = \emptyset$,
$!A \ident !C$, $!B \ident !D \lif !E$
என எடுத்துக்கொள்ள வேண்டும் என்று தெரிகிறது.
எனவே ஒரு நிறுவலின் இறுதி அனுமானம் இதுவாகலாம்:
\[\Axiom$ !C, (!C \land !D) \lif !E \fCenter !D \lif !E$
\RightLabel{\RightR{\lif}}
\UnaryInf$(!C \land !D) \lif !E \fCenter !C \lif (!D \lif !E)$
\DisplayProof\]

% SEGMENT OLP-0702-S02
இந்த எண்ணத்தை மீண்டும் பயன்படுத்தி,
$!D \lif !E$ ஐச் செயலுக்குரிய வாய்பாடாக
எடுத்தால் கிடைப்பது:
\[\Axiom$ !D, !C, (!C \land !D) \lif !E \fCenter !E$
\RightLabel{\RightR{\lif}}
\UnaryInf$!C, (!C \land !D) \lif !E \fCenter !D \lif !E$
\RightLabel{\RightR{\lif}}
\UnaryInf$(!C \land !D) \lif !E \fCenter !C \lif (!D \lif !E)$
\DisplayProof\]
இப்போது $(!C \land !D) \lif !E$ இல்
கவனம் செலுத்தி, \Log{G1c} இன்
$\LeftR{\lif}$ விதியைப் பயன்படுத்த வேண்டும்:
\[\Axiom$ \Gamma \fCenter \Delta, !A$
\Axiom$ !B, \Gamma \fCenter \Delta$
\RightLabel{\LeftR{\lif}}
\BinaryInf$ !A \lif !B, \Gamma \fCenter \Delta$
\DisplayProof\]
இதனால் நமக்குக் கிடைப்பது:
\[\Axiom$!D, !C \fCenter !E, !C \land !D$
\Axiom$!D, !C, !E \fCenter !E$
\RightLabel{\LeftR{\lif}}
\BinaryInf$ !D, !C, (!C \land !D) \lif !E \fCenter !E$
\RightLabel{\RightR{\lif}}
\UnaryInf$!C, (!C \land !D) \lif !E \fCenter !D \lif !E$
\RightLabel{\RightR{\lif}}
\UnaryInf$(!C \land !D) \lif !E \fCenter !C \lif (!D \lif !E)$
\DisplayProof\]

% SEGMENT OLP-0702-S03
மேலுள்ள இடதுபுற உச்சித் தொடரணியில்
உள்ள $!C \land !D$ என்ற !!{formula} வை
முதன்மையாகக் கொண்டு $\RightR{\land}$
விதியைப் பயன்படுத்தினால்:
\[
\Axiom$!D, !C \fCenter !E, !C$
\Axiom$!D, !C \fCenter !E, !D$
\RightLabel{\RightR{\land}}
\BinaryInf$!D, !C \fCenter !E, !C \land !D$
\Axiom$!D, !C, !E \fCenter !E$
\RightLabel{\LeftR{\lif}}
\BinaryInf$ !D, !C, (!C \land !D) \lif !E \fCenter !E$
\RightLabel{\RightR{\lif}}
\UnaryInf$!C, (!C \land !D) \lif !E \fCenter !D \lif !E$
\RightLabel{\RightR{\lif}}
\UnaryInf$(!C \land !D) \lif !E \fCenter !C \lif (!D \lif !E)$
\DisplayProof
\]
இதுவரை நன்றாக உள்ளது. நாம் பயன்படுத்திய
விதிகள் \Log{G1c} மற்றும்~\Log{G3c}
இரண்டிலும் ஒரேவிதமானவை. எனவே இதுவரை
செய்தவை அனைத்தும் \Log{G3c} இலும் பொருந்தும்.
இப்போது கிடைத்துள்ள !!a{proof} துண்டின்
எல்லா உச்சித் தொடரணிகளிலும் ஒரே
!!{formula} இடப்புறத்திலும் வலப்புறத்திலும்
உள்ளது. ஆனால் அவை இன்னும் முற்றிலும்
அடிக்கோள்களாகவில்லை.

% SEGMENT OLP-0702-S04
\Log{G1c} இல் உச்சித் தொடரணிகளுக்கான
!!{proof}s ஐ $!A \Sequent !A$ என்ற வடிவிலான
அடிக்கோள்களிலிருந்து தர வேண்டும்.
தளர்த்தல் விதிகளால் இதை எளிதாகச் செய்யலாம்.
எடுத்துக்காட்டாக, இடதுபுற உச்சித் தொடரணி
$!D, !C \Sequent !E, !C$ ஐப் பின்வருமாறு
!!{prove} செய்யலாம்:
\[
\Axiom$!C \fCenter !C$
\RightLabel{\LeftR{\Weakening}}
\UnaryInf$!D, !C \fCenter !C$
\RightLabel{\RightR{\Weakening}}
\UnaryInf$!D, !C \fCenter !E, !C$
\DisplayProof
\]
மொத்தத்தில் \Log{G1c} இல்
நமக்கு இத்தகைய !!a{proof} உள்ளது:
\[
\Axiom$!C \fCenter !C$
\RightLabel{\LeftR{\Weakening}}
\UnaryInf$!D, !C \fCenter !C$
\RightLabel{\RightR{\Weakening}}
\UnaryInf$!D, !C \fCenter !E, !C$
\Axiom$!D \fCenter !D$
\RightLabel{\LeftR{\Weakening}}
\UnaryInf$!D, !C \fCenter !D$
\RightLabel{\RightR{\Weakening}}
\UnaryInf$!D, !C \fCenter !E, !D$
\RightLabel{\RightR{\land}}
\BinaryInf$!D, !C \fCenter !E, !C \land !D$
\Axiom$!E \fCenter !E$
\RightLabel{\LeftR{\Weakening}}
\UnaryInf$!C, !E \fCenter !E$
\RightLabel{\LeftR{\Weakening}}
\UnaryInf$!D, !C, !E \fCenter !E$
\RightLabel{\LeftR{\lif}}
\BinaryInf$ !D, !C, (!C \land !D) \lif !E \fCenter !E$
\RightLabel{\RightR{\lif}}
\UnaryInf$!C, (!C \land !D) \lif !E \fCenter !D \lif !E$
\RightLabel{\RightR{\lif}}
\UnaryInf$(!C \land !D) \lif !E \fCenter !C \lif (!D \lif !E)$
\DisplayProof
\]
இந்த நிறுவலின் கட்டமைப்பு, குறிப்பாக
அதன் உயரம், $!C$, $!D$, ~$!E$
என்பன எந்த !!{formula}s என்பதனைச்
சார்ந்ததல்ல. மேலுள்ள !!{proof} இல்
இந்தக் குறிகளை எந்த !!{formula}s ஆலும்
மாற்றலாம்; கிடைப்பது \Log{G1c} இல்
மீண்டும் ஓர் !!{proof} ஆகும்.
இந்த \Log{G1c} நிறுவல்
\emph{திட்டவடிவமானது}.

% SEGMENT OLP-0702-S05
\Log{G3c} இல் அடிக்கோள்கள்
$!A, \Gamma \Sequent \Delta, !A$
என்ற வடிவிலான தொடரணிகள்; இங்கு
$!A$ அணு வாய்பாடாக இருக்க வேண்டும்.
எனவே $!D, !C \Sequent !E, !C$
என்பது \Log{G3c} அடிக்கோள் போலத்
தோன்றினாலும்
($\Gamma = \{!D\}$; $\Delta = \{!E\}$),
$!C$ உண்மையில் அணுவாக இருந்தால்தான்
அது \emph{உண்மையான} அடிக்கோள்.
நம் நிறுவல் துண்டு, $!C$, $!D$, $!E$
ஆகியவை அணு !!{formula}s ஆக இருந்தால்தான்
\Log{G3c} இல் முழு நிறுவலாகும்.

முறையாகச் சொன்னால்,
\[\Log{G1c} \Proves (!C \land !D) \lif !E \fCenter !C \lif (!D \lif
!E),\]
என்பதை இன்னும் காட்டவில்லை என்றாலும்,
நடைமுறையில் நாம் செய்ய வேண்டியதும்
செய்யக்கூடியதும் இதுவே. ஏனெனில்
$!A, \Gamma \Sequent \Delta, !A$
என்ற வடிவிலான ஒவ்வொரு தொடரணியும்,
$!A$ அணுவாக இல்லாவிட்டாலும்,
\Log{G3c} இல் நிறுவத்தக்கது.
$!A$ இன் ஆழம்~$\depth{!A}$ மீது
தொகுத்தறிந்து இதைக் காட்டலாம்.

% SEGMENT OLP-0702-S06
\begin{defn}\ollabel{defn:depth}
ஒரு !!a{formula} வின்
\emph{ஆழம்}~$\depth{!A}$
பின்வருமாறு வரையறுக்கப்படுகிறது:
\begin{enumerate}
  \item $!A$ அணுவானால், $\depth{!A} = 0$.
  \item $!A \ident \lnot !B$,
  $!A \ident \lforall[x][!B(x)]$,
  அல்லது $!A \ident \lexists[x][!B(x)]$
  எனில், $\depth{!A} = \depth{!B} + 1$.
  \item $!A \ident (!B \land !C)$,
  $(!B \lor !C)$ அல்லது $(!B \lif
  !C)$ எனில்,
  $\depth{!A} = \max(\depth{!B},\depth{!C}) + 1$.
\end{enumerate}
\end{defn}

% SEGMENT OLP-0702-S07
எந்தச் சொல்~$t$ க்கும்
$\depth{A(x)} = \depth{A(t)}$
என்பதை நிறுவுவது கடினமல்ல.
நமக்கு முக்கியமானது: எந்த விதியிலும்
முதன்மை !!{formula} வின் ஆழம்,
அதன் செயலுக்குரிய !!{formula}s இன்
ஆழத்தைவிட எப்போதும் அதிகம்.
எடுத்துக்காட்டாக:
\begin{align*}
\depth{P(x)} = \depth{P(a)} & = 0,\\
\depth{\lforall[x][P(x)]} & = 1,\\
\depth{(\lforall[x][P(x)] \lif P(a))} & = \max(1,0) + 1 = 2.
\end{align*}
இதைக் கொண்டு $\depth{!A}$ மீது
தொகுத்தறிந்து பின்வரும் முடிவைப் போன்ற
சில முடிவுகளை நிறுவலாம்.

% SEGMENT OLP-0702-S08
\begin{prop}\ollabel{prop:G1c-axioms}
எந்த~$!A$ க்கும்,
$!A \Sequent !A$ என்பதற்கான
$\Log{G1c}$ இல் ஒரு !!a{proof} உள்ளது;
அதில் எல்லா அடிக்கோள்களும் அணுவானவை.
\end{prop}

\begin{proof}
$\depth{!A}$ மீது தொகுத்தறிதல் செய்வோம்.
$\depth{!A} = 0$ எனில் $!A$ அணு வாய்பாடு;
$!A \Sequent !A$ ஏற்கெனவே வேண்டிய
வடிவிலான \Log{G1c} !!{proof} ஆகும்.
$\depth{!A} > 0$ எனில், $!A$ குறைந்த
ஆழமுள்ள !!{formula}s இலிருந்து கட்டப்படுகிறது;
எடுத்துக்காட்டாக $!A \ident \lnot !B$,
$!A \ident (!B \land !C)$ அல்லது
$!A \ident \lforall[x][!B(x)]$.
$\depth{!D} < \depth{!A}$ உடைய எல்லா
$!D$ க்கும் அணு அடிக்கோள்களிலிருந்து
$\Log{G1c} \Proves !D \Sequent !D$
என்பதே தொகுத்தறிதல் கருதுகோள்.

% SEGMENT OLP-0702-S09
$!A \ident \lnot !B$ எனில்,
$\depth{!B} < \depth{!A}$.
தொகுத்தறிதல் கருதுகோளால் அணு
அடிக்கோள்களிலிருந்து $!B \Sequent !B$
என்பதற்கான \Log{G1c} !!a{proof}~$\pi_1$
உள்ளது. இதை $\lnot !B \Sequent \lnot !B$
என்பதற்கான !!a{proof} ஆகப் பின்வருமாறு
நீட்டிக்கலாம்:
\[
\AxiomC{}
\RightLabel{$\pi_1$}
\Deduce$!B \fCenter !B$
\RightLabel{\LeftR{\lnot}}
\UnaryInf$\lnot !B, !B \fCenter $
\RightLabel{\RightR{\lnot}}
\UnaryInf$\lnot !B \fCenter \lnot !B$
\DisplayProof
\]

% SEGMENT OLP-0702-S10
$!A \ident (!B \land !C)$ எனில்,
$\depth{!B} < \depth{!A}$ மற்றும்
$\depth{!C} < \depth{!A}$.
தொகுத்தறிதல் கருதுகோளால்,
அணு அடிக்கோள்களிலிருந்து முறையே
$!B \Sequent !B$,
$!C \Sequent !C$ என்பவற்றுக்கான
\Log{G1c} !!{proof}s~$\pi_1$,
~$\pi_2$ உள்ளன. அவற்றை
$!B \land !C \Sequent !B \land !C$
என்பதற்கான !!a{proof} ஆகப்
பின்வருமாறு நீட்டிக்கலாம்:
\[
\AxiomC{}
\RightLabel{$\pi_1$}
\Deduce$!B \fCenter !B$
\RightLabel{\LeftR{\Weakening}}
\UnaryInf$!B, !C \fCenter !B$
\RightLabel{\LeftR{\land}}
\UnaryInf$!B \land !C, !C\fCenter !B$
\RightLabel{\LeftR{\land}}
\UnaryInf$!B \land !C, !B \land !C \fCenter !B$
\AxiomC{}
\RightLabel{$\pi_2$}
\Deduce$!C \fCenter !C$
\RightLabel{\LeftR{\Weakening}}
\UnaryInf$!B, !C \fCenter !C$
\RightLabel{\LeftR{\land}}
\UnaryInf$!B \land !C, !C \fCenter !C$
\RightLabel{\LeftR{\land}}
\UnaryInf$!B \land !C, !B \land !C \fCenter !C$
\RightLabel{\RightR{\land}}
\BinaryInf$!B \land !C, !B \land !C \fCenter !B \land !C$
\RightLabel{\LeftR{\Contraction}}
\UnaryInf$!B \land !C \fCenter !B \land !C$
\DisplayProof
\]

% SEGMENT OLP-0702-S11
இப்போது $!A \ident \lforall[x][!B(x)]$
என்க. $\lforall[x][!B(x)]$ இல்
வராத !!a{constant}~$c$ ஐ எடுக்கவும்.
அப்போது $\depth{!B(c)} =
\depth{!B(x)} < \depth{!A}$.
தொகுத்தறிதல் கருதுகோளால்,
அணு அடிக்கோள்களிலிருந்து
$!B(c) \Sequent !B(c)$ என்பதற்கான
\Log{G1c} !!a{proof}~$\pi_1$ உள்ளது.
இதை $\lforall[x][!B(x)] \Sequent
\lforall[x][!B(x)]$ என்பதற்கான
!!a{proof} ஆகப் பின்வருமாறு நீட்டிக்கலாம்:
\[
\AxiomC{}
\RightLabel{$\pi_1$}
\Deduce$!B(c) \fCenter !B(c)$
\RightLabel{\LeftR{\lforall}}
\UnaryInf$\lforall[x][!B(x)] \fCenter !B(c)$
\RightLabel{\RightR{\lforall}}
\UnaryInf$\lforall[x][!B(x)] \fCenter \lforall[x][!B(x)]$
\DisplayProof
\]
எஞ்சிய நிலைகள் பயிற்சிகளாக விடப்படுகின்றன.
\end{proof}

% SEGMENT OLP-0702-S12
\begin{prob}
  $!A \ident (!B \lor !C)$,
  $!A \ident (!B \lif !C)$ மற்றும்
  $!A \ident \lexists[x][!B(x)]$
  எனும் நிலைகளைச் செய்து,
  \olref{prop:G1c-axioms} இன்
  நிறுவலை நிறைவு செய்க.
\end{prob}

நமது நோக்கத்திற்கு, $!A \ident \lnot !B$
என்ற நிலையில் பின்வரும் !!{proof} ஐயும்
பயன்படுத்தியிருக்கலாம்:
\[
\AxiomC{}
\RightLabel{$\pi_1$}
\Deduce$!B \fCenter !B$
\RightLabel{\RightR{\lnot}}
\UnaryInf$\fCenter !B, \lnot !B$
\RightLabel{\LeftR{\lnot}}
\UnaryInf$\lnot !B \fCenter \lnot !B$
\DisplayProof
\]
ஆனால் சில தொடரணிக் கணிய முறைமைகளில்
இது செல்லாது. உள்ளுணர்வுவாதத் தொடரணிக்
கணி \Log{G1i}, \Log{G1c} ஐப் போன்றதே;
ஆனால் ஒவ்வொரு தொடரணியின் வலப்புறத்திலும்
அதிகபட்சம் ஒரு !!{formula} மட்டுமே
அனுமதிக்கப்படுகிறது. $\Delta = \emptyset$
எனில் முதல் !!{proof} \Log{G1i} இலும்
!!a{proof} ஆகும்; இரண்டாவது அவ்வாறு ஆகாது.
ஏனெனில் அதன் ஒரு தொடரணியின் வலப்புறத்தில்
$!B$ மற்றும்~$\lnot !B$ இரண்டும் உள்ளன.

% SEGMENT OLP-0702-S13
$!A \ident \lforall[x][!B(x)]$
என்ற நிலையிலோ, \Log{G1c} இலும்
விதிகளை மாற்று வரிசையில் பயன்படுத்த
முடியாது என்பதைக் கவனிக்கவும்.
பின்வருவது \emph{!!a{proof} அல்ல}:
\[
\AxiomC{}
\RightLabel{$\pi_1$}
\Deduce$!B(c) \fCenter !B(c)$
\RightLabel{\RightR{\lforall}}
\UnaryInf$!B(c) \fCenter \lforall[x][!B(x)]$
\RightLabel{\LeftR{\lforall}}
\UnaryInf$\lforall[x][!B(x)] \fCenter \lforall[x][!B(x)]$
\DisplayProof
\]
ஏனெனில் $\RightR{\lforall}$ இன்
தனிமாறி நிபந்தனை மீறப்படுகிறது.

\begin{prop}\ollabel{prop:G3c-axioms}
$!A, \Gamma \Sequent \Delta, !A$
என்ற வடிவிலான எந்தத் தொடரணியும்,
$!A$ அணுவாக \emph{இருக்க வேண்டியதில்லை}
என்றாலும், \Log{G3c} இல் நிறுவத்தக்கது.
\end{prop}

% SEGMENT OLP-0702-S14
\begin{proof}
இந்த நிறுவல் \olref{prop:G1c-axioms}
இன் நிறுவலை மிகவும் ஒத்தது; ஆனால்
தொகுத்தறிதல் கருதுகோளை கவனமாகக்
கையாள வேண்டும். $n = \depth{!A} = 0$
என்ற நிலை மீண்டும் எளிது.
$n=0$ எனில் $!A$ அணு வாய்பாடு;
$!A, \Gamma \Sequent \Delta, !A$
\Log{G3c} இன் அடிக்கோள்.

$n>0$ என்க. மீண்டும் நிலைகளைப்
பிரித்துப் பார்ப்போம்.
தொகுத்தறிதல் கருதுகோள் இதுதான்:
$\depth{!D} < n$ உடைய எல்லா
$!D$ க்கும், எல்லா $\Pi$, $\Lambda$
ஆகியவற்றுக்கும்,
$\Log{G3c} \Proves !D, \Pi \Sequent \Lambda, !D$.
$!A \ident (!B \land !C)$ என்ற நிலையைப்
பார்ப்போம். தொகுத்தறிதல் கருதுகோளை
இருமுறை பயன்படுத்துகிறோம். முதலில்
$!D = !B$, $\Pi = !C, \Gamma$,
$\Lambda = \Delta$ என வைத்து,
$!B, !C, \Gamma \Sequent \Delta, !B$
என்பதற்கான !!a{proof}~$\pi_1$ ஐப்
பெறுகிறோம். அடுத்து
$!D = !C$, $\Pi = !B, \Gamma$,
$\Lambda = \Delta$ என வைத்து,
$!B, !C, \Gamma \Sequent \Delta, !C$
என்பதற்கான !!a{proof}~$\pi_2$ ஐப்
பெறுகிறோம். $!B \land !C, \Gamma \Sequent
\Delta, !B \land !C$ என்பதன்
!!{proof} இதோ:
\[
\AxiomC{}
\RightLabel{$\pi_1$}
\Deduce$!B, !C, \Gamma \fCenter \Delta, !B$
\RightLabel{\LeftR{\land}}
\UnaryInf$!B \land !C, \Gamma \fCenter \Delta, !B$
\AxiomC{}
\RightLabel{$\pi_2$}
\Deduce$!B, !C, \Gamma \fCenter \Delta, !C$
\RightLabel{\LeftR{\land}}
\UnaryInf$!B \land !C, \Gamma \fCenter \Delta, !C$
\RightLabel{\RightR{\land}}
\BinaryInf$!B \land !C, \Gamma \fCenter \Delta, !B \land !C$
\DisplayProof
\]

% SEGMENT OLP-0702-S15
$!A \ident \lforall[x][!B(x)]$
எனில், $\lforall[x][!B(x)]$,
$\Gamma$, ~$\Delta$ ஆகியவற்றில்
வராத !!a{constant} ஒன்றைத் தேர்க.
$!D = !B(c)$,
$\Pi = \lforall[x][!B(x)], \Gamma$
மற்றும் $\Lambda = \Delta$ என வைத்து
தொகுத்தறிதல் கருதுகோளைப் பயன்படுத்துக.
இதனால் $!B(c),
\lforall[x][!B(x)], \Gamma \Sequent \Delta, !B(c)$
என்பதற்கான !!a{proof}~$\pi_1$ கிடைக்கிறது.
அதிலிருந்து பின்வருவது கிடைக்கும்:
\[
\AxiomC{}
\RightLabel{$\pi_1$}
\Deduce$!B(c), \lforall[x][!B(x)], \Gamma \fCenter \Delta, !B(c)$
\RightLabel{\LeftR{\lforall}}
\UnaryInf$\lforall[x][!B(x)], \Gamma \fCenter \Delta, !B(c)$
\RightLabel{\RightR{\lforall}}
\UnaryInf$\lforall[x][!B(x)], \Gamma \fCenter \Delta, \lforall[x][!B(x)]$
\DisplayProof
\]
$c$ ஐத் தேர்ந்த விதத்தால்
$\RightR{\lforall}$ இன் தனிமாறி
நிபந்தனை நிறைவேறுகிறது.
\end{proof}

% SEGMENT OLP-0702-S16
\begin{prob}
  $!A \ident (!B \lif !C)$ மற்றும்
  $!A \ident \lexists[x][!B(x)]$
  எனும் நிலைகளைச் செய்து,
  \olref{prop:G3c-axioms} இன்
  நிறுவலைத் தொடர்க.
\end{prob}

\end{document}

